% !TEX root = ../main.tex
% ============================================================================
%  Informação do relatório — é praticamente só isto que muda entre relatórios
% ============================================================================

\newcommand{\universidade}{Iscte -- Instituto Universitário de Lisboa}
\newcommand{\escola}{Instituto Superior de Ciências do Trabalho e da Empresa}
\newcommand{\curso}{Mestrado em Cibersegurança e Resiliência}
\newcommand{\uc}{Nome da Unidade Curricular}

\newcommand{\titulo}{Título do Relatório}
\newcommand{\subtitulo}{Trabalho Prático n.º 1}
\newcommand{\titulocurto}{Título Curto}   % aparece no cabeçalho das páginas

% Uma linha por autor:  Nome & \email{...} \\
\newcommand{\autores}{%
  Nome Apelido  & \email{nome_apelido@iscte-iul.pt} \\
  Outro Nome    & \email{outro_nome@iscte-iul.pt} \\
}
\newcommand{\grupo}{Grupo 00}             % deixar vazio {} se não aplicável
\newcommand{\docente}{Nome do Docente}

\newcommand{\anoletivo}{Ano Letivo 2026/2027 \textbar{} 1.º Semestre}
\newcommand{\datarelatorio}{\today}

% Logótipo da capa: branco com fundo transparente (vai sobre o azul)
\newcommand{\ficheirologo}{imagens/logo-iscte.png}

% Cor principal (títulos, legendas, capa) — azul Iscte
\newcommand{\corhex}{0E29C2}
